\documentclass[11pt]{article}
\usepackage[T1]{fontenc}
\usepackage{lmodern,microtype,amsmath,amssymb,amsthm,booktabs}
\usepackage[margin=1in]{geometry}
\usepackage[hidelinks]{hyperref}
\newtheorem{lemma}{Lemma}
\newtheorem{corollary}[lemma]{Corollary}
\title{Faster Gaussian resampling and a conditional\\
integer-multiplication saving above $2^{-30}$}
\author{eumemic\thanks{Prepared with assistance from Claude (Anthropic).
This is not independent review or endorsement by Anthropic.}}
\date{October 8, 2026}
\begin{document}
\maketitle
\begin{abstract}
After the compressed complex network, the conditional witness
$\kappa=59/10^{11}$ was capped by the Gaussian resampling maps, whose cost
$O(tp^{3/2+\delta}\alpha)$ per line forces the dimension exponent
$\epsilon<1/5$. We make two changes to those maps. First, a chirp identity
writes every Gaussian sum along a line as one correlation per block, which the
established integer multiplier evaluates in $O(p^{1+\delta})$ per output,
independently of the Gaussian width. Second, a potential-function argument
shows that the powers of the correction $E=N-I$ decay like
$e^{-\pi\alpha^2\sigma n}$ after a burn-in of order $\alpha^2/\theta$ bits,
so the Neumann series needs $O(p/\alpha^2+1/\theta)$ terms rather than
$p/(\alpha^2\theta)$. With $\alpha=\lfloor\sqrt{b/(8d)}\rfloor$ the Gaussian
cost per bit becomes $O(d^2p^\delta)$, and $\epsilon$ may approach $1/2$. The
compressed complex network leaves room for the larger guard parameter this
needs. The result is the conditional witness $\kappa=1479/10^{12}>2^{-30}$, a
factor $1479/590\approx2.51$ above the preceding witness. The underlying
manuscript's interfaces remain assumptions.
\end{abstract}

\section{Scope and dependencies}
The source is OpenAI's \emph{Integer multiplication below $n\log n$}, pinned at
\begin{center}\texttt{adc7f1241b42e322a6451854ab7e4b4c146bf78a},\end{center}
with the compact-control construction and the compressed complex network. The
computational model is a fixed finite alphabet and a fixed number of
one-dimensional tapes. We retain the paired bit network ($\tau=1-296/10^{11}$),
the compressed complex network ($\sigma=1-14/10^9$), the compact-control
movement, reservation, repair and generalized guard, and all transform, error
and rounding interfaces. Only the one-dimensional resampling algorithms, the
Gaussian width $\alpha$ and the assembly parameters change. Numbered results of
Harvey and van der Hoeven~\cite{HarveyHoeven2021} refer to their author-hosted
manuscript.

\section{The one-dimensional interface}
For integers $2\le s<t<2^p$ with $\gcd(s,t)=1$ and $2\le\alpha<\sqrt p$, the
resampling argument of \cite[Section~4]{HarveyHoeven2021} uses the Gaussian
maps $\mathsf S,\mathsf T:\mathbb C^s\to\mathbb C^t$,
\[
 (\mathsf Su)_k=\frac1\alpha\sum_{j\in\mathbb Z}e^{-\pi\alpha^{-2}(j-sk/t)^2}u_j,
 \qquad
 (\mathsf Tu)_k=\sum_{j\in\mathbb Z}e^{-\pi\alpha^2(tj/s-k)^2}u_j,
\]
the nearest indices $q_j=\lfloor tj/s+1/2\rfloor$ with $\beta_j=tj/s-q_j$, the
selection $(Cv)_j=v_{q_j}$, the diagonal $(\mathsf Du)_j=e^{\pi\alpha^2\beta_j^2}u_j$,
and $N=C\mathsf T\mathsf D$. Put $\sigma=t/s$ and $\theta=\sigma-1$. The cited
Theorem~4.2 and Lemma~4.6 give
\begin{equation}\label{eq:gaussian-public-facts}
 \mathsf TP_sF_s=P_tF_t\mathsf S,\qquad
 \|N-I\|<2.01e^{-\pi\alpha^2\theta/2}<2^{-\alpha^2\theta}
 \quad(\alpha^2\theta\ge1).
\end{equation}
The manuscript's no-sort interface keeps the permutations $P_s,P_t$ in the
transform identity and approximates $\mathsf S'=\mathsf S/2$,
$\mathsf J'=N^{-1}/2$ and $\mathsf D'=2^{2-2\alpha^2}\mathsf D$. We change its
hypotheses and time bounds as follows.

\begin{lemma}[One-dimensional interface, revised]
\label{lem:no-sort-resampling}
Assume in addition
\[
 \alpha^{-2}\le\theta\le\tfrac14,\qquad
 n_{\alpha,\theta}:=\min\Bigl\{\Bigl\lceil\frac p{\alpha^2\theta}\Bigr\rceil,
   \Bigl\lceil\frac{p+1}{4\alpha^2}+\frac1\theta\Bigr\rceil\Bigr\}.
\]
Then $P_sF_s=2^{2\alpha^2}B_0P_tF_tA$ with contractions $A=\mathsf S'$ and
$B_0=\mathsf D'\mathsf J'C$, and there are approximations $\widetilde A$ and
$\widetilde B_0$ with errors less than $p^2$, in time
$O((t+p)p^{1+\delta})$ and $O(n_{\alpha,\theta}(t+p)p^{1+\delta})$.
\end{lemma}

\begin{proof}
The exact identity and the norm bounds $\|\mathsf S'\|<3/4$,
$\|\mathsf J'\|<7/8$, $\|\mathsf D'\|<1$ are those of the manuscript: the cited
Lemma~4.6 needs only $\alpha^2\theta\ge1$, which gives $\|N-I\|<0.42$, and
the proof of the cited Proposition~4.7(i) uses only that bound, $\alpha\ge2$ and the cited
Lemma~4.5. The original hypothesis $\theta>p/\alpha^4$ was used only to bound
the cost of the Neumann evaluation. Lemma~4.8 of~\cite{HarveyHoeven2021} gives
$\widetilde{\mathsf D}'$ with error less than $4$ in time $O(tp^{1+\delta})$;
Lemma~\ref{lem:chirped-gaussian} below gives $\widetilde{\mathsf S}'$ and
$\widetilde E$; and the Neumann evaluation of the cited Lemma~4.12, run for
$n_{\alpha,\theta}$ terms with this $\widetilde E$, gives
$\widetilde{\mathsf J}'$ with error less than $3p^2/4$. Its analysis uses only
$\|E\|<1/2$, $\varepsilon(\widetilde E)<p/3$, $n\le p$ and $\|E^n\|\le2^{-p}$,
with $\|E^{n+k}\|\le\|E^n\|\|E\|^k$ for the tail;
Corollary~\ref{cor:neumann-count} supplies the last two. The selection $C$ and
the record adapters cost $O(tp)$ as before.
\end{proof}

\subsection{Faster Gaussian maps}
\label{sec:fast-gaussian}

This subsection supplies the two ingredients used in the proof of
Lemma~\ref{lem:no-sort-resampling}. We keep the notation of that proof and
put $\sigma=t/s$, $\theta=\sigma-1$. Define $q_j$ and $\beta_j$ by the
same formulas for every integer $j$, so that $q_{j+s}=q_j+t$ and
$\beta_{j+s}=\beta_j$. Write $E=N-I$. By \cite[Section~4.2]{HarveyHoeven2021},
\[
 (Eu)_\ell=\sum_{h\ne0}e^{-\pi\alpha^2X_{\ell,h}}u_{\ell+h},\qquad
 X_{\ell,h}=(\sigma h+\beta_\ell)^2-\beta_{\ell+h}^2,
\]
with indices of $u$ taken modulo $s$, and every $X_{\ell,h}$ is
nonnegative.

\begin{lemma}[Powers of the correction]
\label{lem:correction-powers}
Assume $\alpha\geq2$ and $0<\theta<1$. For every integer $n\geq1$,
\[
 \|E^n\|\leq e^{\pi\alpha^2(1/(4\theta)+1/2)}
   \bigl(2.01\,e^{-\pi\alpha^2\sigma}\bigr)^n.
\]
\end{lemma}

\begin{proof}
Expanding the product, $(E^nu)_\ell$ is the sum over
$h_1,\ldots,h_n\in\mathbb Z\setminus\{0\}$ of $e^{-\pi\alpha^2X}u_{j_n}$,
where $j_0=\ell$, $j_i=j_{i-1}+h_i$ and $X=\sum_iX_{j_{i-1},h_i}$.
Moving each subtracted square to the following step,
\[
 X=\sum_{i=1}^nc(j_{i-1},h_i)+\beta_{j_0}^2-\beta_{j_n}^2,\qquad
 c(j,h)=(\sigma h+\beta_j)^2-\beta_j^2=\sigma h(\sigma h+2\beta_j).
\]
Put $\Phi(j)=\sigma(\beta_j^2+1/4)/\theta$. For one step $(j,h)$ let
$y=\beta_j+h\theta$ and $w=q_{j+h}-q_j-h\in\mathbb Z$. Then
$\beta_{j+h}=\beta_j+\sigma h-(q_{j+h}-q_j)=y-w$, so $|y-w|\leq1/2$.
Substitute $\beta_j=y-h\theta$ in $c(j,h)$ and in
$\Phi(j+h)-\Phi(j)=\sigma\bigl((y-w)^2-(y-h\theta)^2\bigr)/\theta$.
Expanding and using $\sigma-\theta=1$ gives
\[
 c(j,h)-\Phi(j+h)+\Phi(j)=\sigma h^2+\frac\sigma\theta\,w(2y-w)
 \geq\sigma h^2,
\]
because $w(2y-w)=w^2+2w(y-w)\geq w^2-|w|\geq0$. Sum over the path.
Since $\sigma/(4\theta)\leq\Phi\leq\sigma/(2\theta)$ and
$\beta_{j_n}^2\leq1/4$,
\[
 X\geq\sigma\sum_ih_i^2-\frac\sigma{4\theta}-\frac14
  =\sigma\sum_ih_i^2-\frac1{4\theta}-\frac12.
\]
Consequently
\[
 \|E^n\|\leq e^{\pi\alpha^2(1/(4\theta)+1/2)}
  \Bigl(\sum_{h\ne0}e^{-\pi\alpha^2\sigma h^2}\Bigr)^n .
\]
As $\pi\alpha^2\sigma>4\pi$, the inner sum is less than
$2e^{-\pi\alpha^2\sigma}(1+2e^{-3\pi\alpha^2\sigma})<2.01e^{-\pi\alpha^2\sigma}$.
\end{proof}

The original estimate $\|E\|^n<2^{-\alpha^2\theta n}$ gains only
$\alpha^2\theta$ bits per term. The lemma gains about
$\pi\alpha^2\log_2e$ bits per term after a burn-in of order
$\alpha^2/\theta$ bits.

\begin{corollary}
\label{cor:neumann-count}
Assume $\alpha\geq2$, $\alpha^2\theta\geq1$ and $\theta<1$, and let
$n_{\alpha,\theta}$ be as in Lemma~\ref{lem:no-sort-resampling}.
Then $n_{\alpha,\theta}\leq p$ and $\|E^{n_{\alpha,\theta}}\|\leq2^{-p}$.
\end{corollary}

\begin{proof}
If $n_{\alpha,\theta}=\lceil p/(\alpha^2\theta)\rceil$, then
$n_{\alpha,\theta}\le p$ and $\|E^{n}\|\le\|E\|^n<2^{-\alpha^2\theta n}\le2^{-p}$
by \eqref{eq:gaussian-public-facts}. Otherwise
$n=n_{\alpha,\theta}\geq(p+1)/(4\alpha^2)+1/\theta$ and $n<\lceil p/(\alpha^2\theta)\rceil\le p$.
Put $\kappa_0=\pi\log_2e$, so $4.531<\kappa_0<4.534$, and note
$\log_22.01<1.01$. Taking base-two logarithms in
Lemma~\ref{lem:correction-powers},
\[
 \log_2\|E^n\|\leq\kappa_0\alpha^2\Bigl(\frac1{4\theta}+\frac12\Bigr)
   -n(\kappa_0\alpha^2\sigma-1.01).
\]
Since $\sigma>1$ and $\alpha\ge2$, the last bracket is at least
$4\alpha^2$. Also $4.534(1/(4\theta)+1/2)\le4/\theta$ because
$\theta<1$. Hence
$\log_2\|E^n\|\leq4\alpha^2/\theta-4\alpha^2n\leq-p-1$.
\end{proof}

\begin{lemma}[Gaussian maps by chirped correlation]
\label{lem:chirped-gaussian}
Assume the hypotheses of Lemma~\ref{lem:no-sort-resampling}. There are
fixed multitape algorithms computing approximations
$\widetilde{\mathsf S}':\mathbb D_p^s\to\mathbb D_p^t$ of $\mathsf S'$
and $\widetilde E:\mathbb D_p^s\to\mathbb D_p^s$ of $E$, with errors less
than $16p$ and $p/3$ respectively, outputs in the disk grid, and time
$O((t+p)p^{1+\delta})$ each. Inputs and outputs occur in increasing
coordinate order. The constants are uniform in $p,s,t,\alpha$.
\end{lemma}

\begin{proof}
\emph{Chirp identity.} For a rational $y$ and integers $a,b$, put
$\hat a=a+y/\sigma$. Expanding both sides and using $\sigma-\theta=1$,
\[
 (\sigma\hat a-b)^2
   =\sigma\theta\hat a^2+\sigma(\hat a-b)^2-\theta b^2 .
\]
Hence, for rational $\lambda>0$,
\begin{equation}\label{eq:chirp-split}
 e^{-\pi\lambda(\sigma a+y-b)^2}
  =e^{\pi\lambda\theta b^2}\,
   e^{-\pi\lambda\sigma\theta(a+y/\sigma)^2}\,
   e^{-\pi\lambda\sigma(a-b+y/\sigma)^2}.
\end{equation}
For inputs $f_a$, $0\le a<L_A$, the sums
$F_b=\sum_ae^{-\pi\lambda(\sigma a+y-b)^2}f_a$, $0\le b<L_B$, are therefore
$F_b=e^{\pi\lambda\theta b^2}c_b$, where
\[
 c_b=\sum_ag_aK_{a-b},\qquad
 g_a=e^{-\pi\lambda\sigma\theta(a+y/\sigma)^2}f_a,\qquad
 K_r=e^{-\pi\lambda\sigma(r+y/\sigma)^2}.
\]
This is one correlation of finite sequences; $K_r\in(0,1]$ and
$|g_a|\le|f_a|$.

\emph{One block.} Let $P\ge p$ be an integer, $|f_a|\le F$ with $F\ge1$,
and $e^{\pi\lambda\theta b^2}\le2^{\Gamma}$ for $0\le b<L_B$, with
$\Gamma\le2P$. Suppose approximations $\widetilde f_a$ with
$|\widetilde f_a-f_a|\le2^{4-P}F$ are available. All exponents below are $\pi$
times rationals whose numerators and denominators have $O(p)$ bits.
Lemmas~2.13 and~2.14 of \cite{HarveyHoeven2021}, at precision $P$, give the
decaying factors and kernel values, and the numbers
$2^{-\Gamma}e^{\pi\lambda\theta b^2}$, each with absolute error at most
$2^{1-P}$. Multiply each $\widetilde f_a$ by its factor and round to
precision $P$.
Represent every resulting real or imaginary part as an integer multiple of
$2^{-P}$, split it into its nonnegative and negative parts, and pack each
sequence into one integer with slots of $2P+\lceil\log_2(FL_A)\rceil+8$
bits. A constant number of products by the established multiplier
\cite[Theorem~1.1]{HarveyHoeven2021} then gives every rounded $c_b$
exactly, because no slot overflows. Multiply by the approximate factor
$2^{-\Gamma}e^{\pi\lambda\theta b^2}$, shift by $\Gamma$, and round to
precision $p$.

The rounded $g_a$ and $K_r$ have errors at most $2^{5-P}F$ and $2^{1-P}$,
so the rounded correlation differs from $c_b$ by at most $2^{7-P}FL_A$,
while $|c_b|\le FL_A$. After the shift, the final factor has absolute
error at most $2^{\Gamma+1-P}$. Therefore the computed value differs from
$F_b$ by at most $2^{\Gamma+9-P}FL_A$, which is at most $2^{-p-2}$ whenever
\[
 P\ge p+\Gamma+\lceil\log_2(FL_A)\rceil+11.
\]
The block uses $O(L_A+L_B)$ exponentials and products at precision $P$
and integer products on $O((L_A+L_B)(P+\log(FL_A)))$ bits. When
$P=O(p)$ and $\log(FL_A)=O(p)$, its cost is
$O((L_A+L_B)p^{1+\delta})$.

\emph{The map $\mathsf S'$.} Here
$(\mathsf S'u)_k=(2\alpha)^{-1}\sum_je^{-\pi\lambda(\sigma j-k)^2}u_j$
with $\lambda=(\sigma\alpha)^{-2}$. Let $m=\lceil\sqrt p\rceil\alpha$, the
window of \cite[Lemma~4.9]{HarveyHoeven2021}. Partition $0\le k<t$ into
consecutive blocks of at most $m$ outputs. For a block starting at $k_0$,
take inputs $j=j_0+a$, $0\le a<L_A=3m+3$, where $j_0=\lfloor k_0/\sigma\rfloor-m$,
with the periodic indices of $u$, and put $k=k_0+b$, $y=\sigma j_0-k_0$.
These inputs contain every $j$ with $|j-k/\sigma|<m$ for every $k$ in the
block. The block sums therefore omit only terms omitted by the truncated sums
of that lemma, so they differ from $(\mathsf S'u)_k$ by less than $2^{-p}$.
Here $f_a=u_{j_0+a}$ is exact, $F=1$, and for $b<m$,
$\pi\lambda\theta b^2\le\pi\theta(\sqrt p+1)^2\le(\pi/4)(1.21p)$ because
$\theta\le1/4$ and $p>100$, so $\Gamma=\lceil1.4p\rceil$ suffices. Take
$P=3p$, and apply the factor $(2\alpha)^{-1}$ by one more product at
precision $P$. Each output then has error less than
$2^{-p}(1+1/4+1)<3\cdot2^{-p}$, and its norm is less than
$\|\mathsf S'\|+3\cdot2^{-p}<1$.

\emph{The map $E$.} With $f_j=e^{\pi\alpha^2\beta_j^2}u_j$ and
$\lambda=\alpha^2$, the definition of $N$ gives
\[
 (Nu)_\ell=\sum_je^{-\pi\lambda(\sigma j-q_\ell)^2}f_j,
\]
and its term $j=\ell$ is exactly $u_\ell$. The inputs $f_j$ are formed by
Lemma~2.14 of \cite{HarveyHoeven2021} for
$2^{-\lceil1.14\alpha^2\rceil}e^{\pi\alpha^2\beta_j^2}$, a shift, and one
product with $u_j$; their errors are at most $2^{4-P}F$. Let
$m=\lceil\sqrt p/(2\alpha)\rceil$, the window of
\cite[Lemma~4.11]{HarveyHoeven2021}. Partition $0\le\ell<s$ into blocks
of at most $m$ indices. For a block starting at $\ell_0$, take inputs
$j=j_0+a$, $0\le a<L_A=3m+1$, with $j_0=\ell_0-m$, and evaluate the
correlation at every integer $b=q_{\ell_0}+b'$, $0\le b'<L_B=2m+2$; these
include all $q_\ell$ of the block because $\sigma<2$. Put
$y=\sigma j_0-q_{\ell_0}$. Then $F=e^{\pi\alpha^2/4}<2^{1.14\alpha^2}\le2^{1.14p}$
and
\[
 \pi\lambda\theta b'^2\le\pi\theta\alpha^2(2m+2)^2
   \le\frac\pi4\bigl(\sqrt p+4\alpha\bigr)^2<\frac{25\pi}4p,
\]
so $\Gamma=\lceil29p\rceil$ and $P=34p$ suffice, using
$(\sqrt p+4\alpha)^2<25p$ and $\theta\le1/4$. Select the outputs
$b=q_\ell$ in order, subtract the exact $u_\ell$, and round. The block
contains every term $0<|h|\le m$ of the truncated sum in that lemma, so
the computed $(\widetilde Eu)_\ell$ differs from $(Eu)_\ell$ by less than
$(3+1+1)2^{-p}<(p/3)2^{-p}$. Its norm is less than
$\|E\|+5\cdot2^{-p}<1$.

\emph{Cost and tape movement.} The selected indices $q_\ell$, the numerators
$s\beta_\ell$, and the block offsets are generated by the
counters already used for $C$, in $O(p)$ time per index. Before a line,
copy $u$ with $3m+3$ periodic records appended at each end, in $O((s+m)p)$
time; then every block reads one consecutive interval. Consecutive input
intervals overlap in $O(m)$ records, so all reads, packing and unpacking
cost $O((t+m)p)$. There are $O(t/m+1)$ blocks, each costing
$O(mp^{1+\delta})$; in both cases $m<2p$. This gives the time bound
$O((t+p)p^{1+\delta})$.
\end{proof}


\section{Parameters and assembly}
The tensor interface applies the one-dimensional maps line by line. With
$n_*=\max_in_{\alpha,\theta_i}$ and every $t_i\ge p$, its time becomes
$\sum_i(T/t_i)O(n_*(t_i+p)p^{1+\delta})=O(dTp^{1+\delta}(1+n_*))$.

\paragraph{Gaussian width.} With $b=\lceil\log_2n\rceil$, $p=6b$,
$d=\lfloor b^\epsilon\rfloor$ and $\eta=1/(4d)$ as in the manuscript, set
\[
 \alpha=\Bigl\lfloor\sqrt{b/(8d)}\Bigr\rfloor,\qquad\gamma=2d\alpha^2\le b/4.
\]
The bound $\gamma\le b/4$ is the one used by the retained final rounding
argument. For large $n$, $b\ge96d$, so $\alpha\ge2$ and $\alpha^2\ge b/(16d)$;
also $\alpha^2<p$. The chosen primes satisfy
$\eta<\theta_i<1/(2d-1)\le1/4$, so $\alpha^2\theta_i>b/(64d^2)\ge1$ once
$b\ge64d^2$, which holds eventually because $2\epsilon<1$. Hence
\[
 n_{\alpha,\theta_i}\le\frac{p+1}{4\alpha^2}+\frac1{\theta_i}+1
 \le\frac{(6b+1)16d}{4b}+4d+1\le30d.
\]
The original choice $\alpha=\lceil(32db)^{1/4}\rceil$ would need
$\lceil p/(\alpha^2\theta_i)\rceil$ Neumann terms, of order $(dp)^{1/2}$;
with the larger width the old windowed Gaussian sums would cost
$\Theta(p/\sqrt d)$ times more per output than the chirped correlations.

\paragraph{Cost row.} Each axis has at most $T/t_i$ lines of length
$t_i>p$, so the Gaussian maps cost $O(d^2Tp^{1+\delta})$, that is
$O(d^2p^\delta)$ per bit, with power $2\epsilon+\delta$. This replaces the
power $3/4+\delta+5\epsilon/4$. All other rows of the cost table are unchanged.

\paragraph{Parameters.} Put $a=296/10^{11}$, $a_{\rm c}=14/10^9$ and choose
\begin{gather*}
 \tau=1-a,\quad\sigma=1-a_{\rm c},\quad
 \epsilon=\frac{49999}{100000},\quad c=\frac{9999}{10000},\quad
 \beta=\frac{19}{25},\quad\zeta=\frac1{10000},\quad\delta=\frac1{10^6},\\
 C_1=5-4\beta+\zeta=\frac{19601}{10000},\quad
 \lambda=1-\frac{29595}{10^{13}},\quad
 \lambda'=1-\frac{2959}{10^{12}},\quad
 \kappa=\frac{1479}{10^{12}}.
\end{gather*}
Since $\sigma<\tau$, the compact-control internal exponent is $\chi=\tau$.
Exact comparisons give
\[
 \max\{\tau,\sigma,\chi\}<\lambda<\lambda'<1,\quad
 \max\{\sigma+\beta(1-\sigma),\,1-c\}<\lambda',\quad
 \epsilon C_1=\frac{980030399}{10^{9}}<1.
\]
At this $\epsilon$ the guard condition $\epsilon C_1<1$ forces $\beta>3/4$,
and the leaf condition forces $(1-\beta)a_{\rm c}>1-\lambda'$. Both hold only because the
compressed complex saving $a_{\rm c}$ exceeds $a$ by a factor $4.7$; with the
original complex saving $418/10^{12}$ no $\beta$ satisfies both. The remaining
setup conditions $\epsilon(1+c)<1$, $1-2\epsilon>0$ and $\epsilon+\delta<1$
hold.

The seven assembly margins are
\begin{align*}
 g_1&=1-\epsilon(1+c),&g_2&=\epsilon c(1-\tau),\\
 g_3&=\epsilon(1-\lambda'),&g_4&=(1-\tau)(1-\epsilon),\\
 g_5&=1-\delta-2\epsilon,&g_6&=1-\delta-\epsilon,\qquad g_7=\epsilon,
\end{align*}
and exact arithmetic gives
\[
 G_*=\min_ig_i=g_3=\frac{147947041}{10^{17}},\qquad
 G_*-\kappa=\frac{47041}{10^{17}}>0.
\]
This strict gap absorbs all remaining fixed powers of $\log p$, and the
retained assembly yields
\[
 T(n)=O\bigl(n(\log n)^{1-1479/10^{12}}\bigr),\qquad
 2^{-30}<\frac{1479}{10^{12}}<2^{-29},
\]
conditional on the identified upstream and extension proofs.

\paragraph{The next ceiling.} Now $g_2,g_3,g_4\le a\min\{\epsilon,1-\epsilon\}$,
so these networks and assembly inequalities give $\kappa<a/2<2^{-29}$. The
witness is above $99.9\%$ of this bound. Further progress needs a stronger
bit network; the complex network would then need to keep
$(1-\beta)a_{\rm c}>a$ with $\beta>3/4$.

\section{Reproducibility and review boundary}
\path{scripts/fast_gaussian.py} checks the witness, the guard, the exact
step inequality of Lemma~\ref{lem:correction-powers} on several $(s,t)$, and the
Neumann counts; its certificate is \path{certificates/fast-gaussian.json}.
\path{tests/test_fast_gaussian.py} verifies the chirp identity exactly, the
power bound against explicit $E$ matrices at $160$ decimal digits, the block
correlation against direct sums, negative parameter cases, and the patch
text. \path{scripts/make_fast_gaussian_patch.py} writes
\path{patches/fast-gaussian-30.patch}. These checks support the written
arguments; they are not a formal verification of the upstream theorem, and
the new construction has not received independent review.

\begin{thebibliography}{1}
\bibitem{HarveyHoeven2021}
D.~Harvey and J.~van der Hoeven, \emph{Integer multiplication in time
$O(n\log n)$}, Ann. of Math. (2) \textbf{193} (2021), no.~2, 563--617.
\end{thebibliography}
\end{document}
